\section{Chebyshev polynomials for spectral filters}

\begin{theorem}[Exponential growth on the positive half-line]
    \label{thm:chebyshev_half_exp_sqrt}
    \lean{Polynomial.Chebyshev.half_exp_sqrt_le_eval_T}\leanok
    For every natural number $n$ and every real $x\ge1$, the Chebyshev
    polynomial of the first kind satisfies
    \begin{align}
        T_n(x) & \ge\frac12\exp\left(2n\sqrt{\frac{x-1}{x+1}}\right).
        \label{eq:chebyshev_half_exp_sqrt}
    \end{align}
    This is the scalar estimate used in
    \cite{Arad2013AreaLaw,Arad2016RigorousRG}.
\end{theorem}
\begin{proof}\leanok
    Put $r=\sqrt{(x-1)/(x+1)}$. Then $0\le r<1$ and
    $\cosh(\log((1+r)/(1-r)))=x$. The logarithmic inequality
    $r\le\tfrac12\log((1+r)/(1-r))$ gives $2r\le\operatorname{arcosh}x$.
    Therefore
    \begin{align}
        T_n(x) & =\cosh(n\operatorname{arcosh}x)
        \ge\tfrac12\exp(n\operatorname{arcosh}x)
        \ge\tfrac12\exp(2nr).\notag
    \end{align}
    These identities also apply when $n=0$ or $x=1$.
\end{proof}


\begin{theorem}[Normalized polynomial on an isolated spectral interval]
    \label{thm:chebyshev_interval_polynomial}
    \lean{Polynomial.Chebyshev.exists_interval_shrinking_polynomial}\leanok
    \uses{thm:chebyshev_half_exp_sqrt}
    Let $a<b\le u$ be real numbers and let $n\in\mathbb N$.
    There is a real polynomial $p$ of degree exactly $n$ with $p(a)=1$.
    It satisfies $p(x)\ge1$ for every $x\le a$, and for every $x\in[b,u]$,
    \begin{align}
        |p(x)| & \le2\exp\left(-2n\sqrt{\frac{b-a}{u-a}}\right).
        \label{eq:chebyshev_interval_polynomial}
    \end{align}
    This is the scalar polynomial construction in
    \cite{Arad2013AreaLaw,Arad2016RigorousRG}, including degree zero and
    a singleton upper interval.
\end{theorem}
\begin{proof}\leanok
    \uses{thm:chebyshev_half_exp_sqrt}
    First suppose $b<u$. Put
    \begin{align}
        f(x) & =\frac{u+b-2x}{u-b},\notag          \\
        p(x) & =\frac{T_n(f(x))}{T_n(f(a))}.\notag
    \end{align}
    The affine map $f$ takes $[b,u]$ into $[-1,1]$, while $f(a)>1$.
    Thus $|T_n(f(x))|\le1$ on $[b,u]$ and $T_n(f(a))\ge1$.
    Theorem~\ref{thm:chebyshev_half_exp_sqrt} and
    $(f(a)-1)/(f(a)+1)=(b-a)/(u-a)$ give
    (\ref{eq:chebyshev_interval_polynomial}).
    The denominator is positive, so the polynomial has degree $n$ and
    takes value one at $a$. Monotonicity of $\cosh(n\operatorname{arcosh}y)$
    for $y\ge1$, together with the decreasing affine map $f$, gives
    $p(x)\ge1$ for $x\le a$.

    If $b=u$, use $p(x)=((u-x)/(u-a))^n$.
    It has degree $n$, takes value one at $a$ and is at least one to the left.
    For $n>0$ it vanishes on the upper interval $\{u\}$; for $n=0$ it is
    the constant one and the required upper bound is two.
\end{proof}


\begin{theorem}[Hermitian filtering of low and high spectral subspaces]
    \label{thm:chebyshev_spectral_filter}
    \lean{Matrix.IsHermitian.exists_polynomial_spectral_filter}\leanok
    \uses{thm:chebyshev_interval_polynomial}
    Let $H$ be a finite Hermitian matrix whose real spectrum is bounded above
    by $u$. Let $a<b\le u$ and $n\in\mathbb N$, and put
    $\delta=2\exp(-2n\sqrt{(b-a)/(u-a)})$.
    There is one polynomial $p$ satisfying
    Theorem~\ref{thm:chebyshev_interval_polynomial} such that $K=p(H)$
    is Hermitian, commutes with $H$, and obeys
    \begin{align}
        P_-KP_-  & \ge P_-,\notag   \\
        \|KP_+\| & \le\delta,\notag
    \end{align}
    where $P_-=\mathbf1_{(-\infty,a]}(H)$ and
    $P_+=\mathbf1_{[b,\infty)}(H)$.
    These are the spectral estimates used in
    \cite{Arad2013AreaLaw,Arad2016RigorousRG}. The additional source
    positivity clause is corrected in
    \cite{gap:alvv16_chebyshev_filter_positivity}.
\end{theorem}
\begin{proof}\leanok
    \uses{thm:chebyshev_interval_polynomial}
    Choose the scalar polynomial once. Functional calculus preserves its
    Hermiticity and commutation with $H$.
    The function $\mathbf1_{t\le a}p(t)\mathbf1_{t\le a}-\mathbf1_{t\le a}$
    is nonnegative, giving the low-energy order inequality.
    On the spectrum, $|p(t)\mathbf1_{b\le t}|\le\delta$, since retained
    spectral values belong to $[b,u]$. Isometric functional calculus gives
    the operator norm bound. Every function on the finite real spectrum
    is continuous there, including these threshold functions.
\end{proof}

\begin{theorem}[Approximation of an isolated-energy projection]
    \label{thm:chebyshev_isolated_energy_projection}
    \lean{Matrix.IsHermitian.exists_polynomial_approximating_isolated_energy_projection}
    \leanok
    \uses{thm:chebyshev_spectral_filter}
    Let $H$ be a finite Hermitian matrix with spectrum in $\{a\}\cup[b,u]$,
    where $a<b\le u$, and let $n\in\mathbb N$.
    There is a polynomial $p$ of degree exactly $n$, normalized by $p(a)=1$,
    such that for $K=p(H)$ and the spectral projection $P$ at $a$,
    \begin{align}
        KP      & =P=PK,\notag                                          \\
        \|K-P\| & \le2\exp\left(-2n\sqrt{\frac{b-a}{u-a}}\right).\notag
    \end{align}
    The multiplicity of $a$ is unrestricted; if $a$ is absent, $P=0$.
\end{theorem}
\begin{proof}\leanok
    \uses{thm:chebyshev_spectral_filter}
    Use the same polynomial as in Theorem~\ref{thm:chebyshev_spectral_filter}.
    The spectral assumption gives $P=1-P_+$, and normalization at $a$
    gives $KP=P=PK$. Hence $K-P=KP_+$, whose norm has the required bound.
\end{proof}

\begin{theorem}[A ground-state Chebyshev filter from a global gap]
    \label{thm:chebyshev_ground_filter}
    \lean{Matrix.IsHermitian.exists_chebyshev_ground_filter}
    \leanok
    Let $H$ be a finite Hermitian matrix, and let $\Gamma$ be a unit vector
    satisfying $H\Gamma=E_0\Gamma$ and
    \begin{align*}
        H-E_0\mathbb 1\ge\gamma(\mathbb 1-|\Gamma\rangle\langle\Gamma|),
        \qquad \gamma>0.
    \end{align*}
    Suppose $E_0+\gamma\le u$ and every spectral value of $H$ is at most $u$.
    For every $\ell\ge0$ there is a real polynomial $p$ of degree exactly $\ell$
    with $p(E_0)=1$. Put
    \begin{align*}
        K=p(H),\qquad P_\Gamma=|\Gamma\rangle\langle\Gamma|,\qquad
        A_\ell=2\exp\left(-2\ell\sqrt{\frac{\gamma}{u-E_0}}\right).
    \end{align*}
    Then $K\Gamma=\Gamma$ and $\|K-P_\Gamma\|\le A_\ell$.
    Every unit vector $\psi\perp\Gamma$ satisfies
    $K\psi\perp\Gamma$ and $\|K\psi\|^2\le A_\ell^2$.
    These are the ground-state and shrinking conditions used in
    \cite[Section~4, Lemma~4.1]{Arad2013AreaLaw}.
\end{theorem}
\begin{proof}
    \leanok
    \uses{thm:chebyshev_isolated_energy_projection,thm:global_gap_ground_projection}
    The global gap gives the spectral alternative
    $E_0$ or $[E_0+\gamma,u]$.
    Construct the normalized polynomial by the isolated-energy filter theorem.
    Its canonical low-energy projection equals $P_\Gamma$ by the global gap,
    and $KP_\Gamma=P_\Gamma K=P_\Gamma$.
    Since $P_\Gamma\Gamma=\Gamma$, the filter fixes $\Gamma$.
    If $\psi\perp\Gamma$, then $P_\Gamma\psi=0$ and
    $P_\Gamma K\psi=0$, proving invariance of the orthogonal complement.
    Finally,
    $\|K\psi\|=\|(K-P_\Gamma)\psi\|\le\|K-P_\Gamma\|\|\psi\|\le A_\ell$;
    squaring gives the shrinking coefficient.
\end{proof}
